\scanpage{171}
\begin{center}- 43 -\end{center}\hfill IX

-\uline{donc on a pour l'action de} $g$ \uline{sur des fonctions} $\varphi$ :
\[
 g^*(\lambda\varphi)=f^*(\lambda)g^*(\varphi)
 =\lambda^{q^{-1}}g^*(\varphi)=\lambda^{p^{-r}}g^*(\varphi)
 \qquad(\lambda\in k)\qquad.
\]
\uline{Considérons l'automorphisme} $g_\ell$ \uline{induit par} $g$ \uline{sur}
$T_\ell(A(k))$ \uline{par ``transport de structure''}, \uline{où} $\ell$
\uline{est un nombre premier} $\neq p$. \uline{Alors la trace de} $g_\ell$
(\uline{plus généralement}, \uline{tout coefficient de son polynôme caractéristique})
\uline{est un entier rationnel}, \uline{indépendant de} $\ell$, \uline{et ses valeurs propres
sont des entiers algébriques de valeur absolue égale à} $q^{1/2}$ \uline{ou} $q$.
\uline{Si} $A$ \uline{est un schéma abélien}, \uline{les valeurs absolues sont} $q^{1/2}$ ;
\uline{si} $A$ \uline{est affine}, \uline{les valeurs propres sont de la forme} $q\epsilon_i$,
\uline{où les} $\epsilon_i$ \uline{sont des racines de l'unité}.

En effet, considérons l'endomorphisme $f_q^A$ ``puissance q.ième''
du schéma absolu sous-jacent à $A$, qui est défini comme l'identité
sur les points et l'élévation à la puissance q.ième sur le faisceau
structural. C'est un endomorphisme du schéma en groupes $A$ compatible
avec l'endomorphisme $f_q^{\operatorname{Spec}(k)}=f^{-1}$ sur la base $\operatorname{Spec}(k)$, donc le
composé
\[
 (4.3.3)\qquad F=f_q^A\,g
\]
est un endomorphisme du schéma en groupes $A$, compatible avec l'identité
sur $\operatorname{Spec}(k)$, i.e. c'est un k-endomorphisme du schéma en groupes $A$.
Je dis que l'automorphisme $\Gamma_g$ de $A(k)$ défini par $g$ par transport de
structure, explicité par la formule
\[
 (4.3.4)\qquad \Gamma_g(u)=g\circ u\circ f^{-1}\qquad\text{pour }u\in A(k)\qquad,
\]
peut aussi de décrire par
\[
 (4.3.5)\qquad\Gamma_g(u)=F(u)\qquad\text{pour }u\in A(k)\qquad.
\]
\margin{361}
\margin{361}
\margin{170}
